\chapter{图表与公式排版}
\label{chap:layout}
\section{图的编号与标题}
如图~\ref{fig:workflow}~所示，研究过程可表示为资料收集、方法设计和结果分析三个阶段。
图题置于图的下方，图注置于图题上方。硕士论文可以使用中文图题；本例同时展示可选的双语图题，中文在上。
\begin{figure}[htbp]
  \centering
  \begin{minipage}{0.85\textwidth}
    \centering
    \setlength{\fboxsep}{10pt}
    \fbox{资料收集}\quad$\longrightarrow$\quad
    \fbox{方法设计}\quad$\longrightarrow$\quad\fbox{结果分析}
    \medskip
    \gzunote{注：方框为排版演示，不表示特定实验步骤。}
    \gzubicaption{研究流程示意}{Schematic of the research workflow}
    \label{fig:workflow}
  \end{minipage}
\end{figure}

\section{三线表}
表~\ref{tab:parameters}~展示三线表的写法。表题位于表格上方，表格内容为小五号宋体，顶线和底线宽度为 1.5 磅，中线宽度为 0.75 磅。
\begin{table}[htbp]
  \centering
  \gzubicaption{示例参数及其含义}{Example parameters and their meanings}
  \label{tab:parameters}
  \begin{tabular}{clc}
    \toprule
    符号 & 含义 & 示例值 \\
    \midrule
    $n$ & 样本数量 & 120 \\
    $d$ & 特征维数 & 8 \\
    $\alpha$ & 权重系数 & 0.5 \\
    \bottomrule
  \end{tabular}
  \gzunote{注：数值仅用于展示表格排版。}
\end{table}

\section{数学公式}
如式~\eqref{eq:mean}~所示，算术平均值由所有观测值之和除以样本数量得到。
公式单独成行，编号按章排列并在右侧对齐。
\begin{equation}
  \bar{x}=\frac{1}{n}\sum_{i=1}^{n}x_i .
  \label{eq:mean}
\end{equation}
\gzuequationwhere{$\bar{x}$ 为算术平均值；$n$ 为样本数量；$x_i$ 为第 $i$ 个观测值。}

\subsection{多行公式}
较长公式可使用对齐环境，在关系符号附近换行。式~\eqref{eq:variance}~给出一个两行公式示例。
\begin{equation}
  \begin{aligned}
    s^2 &= \frac{1}{n-1}\sum_{i=1}^{n}(x_i-\bar{x})^2 \\
        &= \frac{1}{n-1}\left(\sum_{i=1}^{n}x_i^2-n\bar{x}^{\,2}\right).
  \end{aligned}
  \label{eq:variance}
\end{equation}

\section{跨页表格}
表~\ref{tab:long}~故意设置为较长的表格，以检查跨页时的表号、续表标记和重复表头。
跨页表格应直接使用 longtable 环境，不要嵌套在 table 浮动体中。该示例采用硕士论文允许的中文表题。
\begin{longtable}{cll}
  \caption{跨页数据排版示例}\label{tab:long}\\
  \toprule 编号 & 项目 & 说明 \\ \midrule
  \endfirsthead
  \caption[]{跨页数据排版示例（续）}\\
  \toprule 编号 & 项目 & 说明 \\ \midrule
  \endhead
  \bottomrule
  \endfoot
  1 & 项目一 & 示例数据 \\
  2 & 项目二 & 示例数据 \\
  3 & 项目三 & 示例数据 \\
  4 & 项目四 & 示例数据 \\
  5 & 项目五 & 示例数据 \\
  6 & 项目六 & 示例数据 \\
  7 & 项目七 & 示例数据 \\
  8 & 项目八 & 示例数据 \\
  9 & 项目九 & 示例数据 \\
  10 & 项目十 & 示例数据 \\
  11 & 项目十一 & 示例数据 \\
  12 & 项目十二 & 示例数据 \\
  13 & 项目十三 & 示例数据 \\
  14 & 项目十四 & 示例数据 \\
  15 & 项目十五 & 示例数据 \\
  16 & 项目十六 & 示例数据 \\
  17 & 项目十七 & 示例数据 \\
  18 & 项目十八 & 示例数据 \\
  19 & 项目十九 & 示例数据 \\
  20 & 项目二十 & 示例数据 \\
  21 & 项目二十一 & 示例数据 \\
  22 & 项目二十二 & 示例数据 \\
  23 & 项目二十三 & 示例数据 \\
  24 & 项目二十四 & 示例数据 \\
  25 & 项目二十五 & 示例数据 \\
  26 & 项目二十六 & 示例数据 \\
  27 & 项目二十七 & 示例数据 \\
  28 & 项目二十八 & 示例数据 \\
  29 & 项目二十九 & 示例数据 \\
  30 & 项目三十 & 示例数据 \\
  31 & 项目三十一 & 示例数据 \\
  32 & 项目三十二 & 示例数据 \\
  33 & 项目三十三 & 示例数据 \\
  34 & 项目三十四 & 示例数据 \\
  35 & 项目三十五 & 示例数据 \\
  36 & 项目三十六 & 示例数据 \\
  37 & 项目三十七 & 示例数据 \\
  38 & 项目三十八 & 示例数据 \\
  39 & 项目三十九 & 示例数据 \\
  40 & 项目四十 & 示例数据 \\
  41 & 项目四十一 & 示例数据 \\
  42 & 项目四十二 & 示例数据 \\
  43 & 项目四十三 & 示例数据 \\
  44 & 项目四十四 & 示例数据 \\
  45 & 项目四十五 & 示例数据 \\
\end{longtable}
